\section{Research Questions}
\label{sec:questions}

\subsection{Original aim and questions}
The proposal's aim was to adapt and evaluate a spatially aware VAE
(Spatial-msiPL) capable of autonomously identifying spatially coherent
molecular peaks in untargeted MSI data, as a computationally efficient
alternative to 3D convolutional networks. It asked three questions:
\begin{enumerate}[label=P\arabic*.]
  \item How can localised spatial covariance be integrated into the
        probabilistic encoding process of a standard 1D VAE without the
        computational and memory overhead associated with 3D convolutional
        networks?
  \item To what extent does an explicit spatial-coherence penalty, combined
        with physical noise augmentation, isolate true biomolecular signals
        from stochastic matrix noise, compared with spatially blind baselines?
  \item To what extent does this spatially regularised framework improve the
        identification of biologically relevant, spatially structured peaks
        compared with spatially blind models such as msiPL?
\end{enumerate}

\subsection{Questions answered in this report}
Development work showed that the penalty and augmentation did not pass their
gates (\cref{tab:proposal}). Fair scoring then showed that the answer to P3
depended more on the ranking step than on spatial input. The report therefore
answers four questions that separate the ingredients. Each is answered with
paired comparisons across sections under a rule fixed before the results
(\cref{sec:design}).
\begin{enumerate}[label=RQ\arabic*.]
  \item \textbf{Baseline.} Does the label-free VAE--GMM pipeline select peaks
        that agree better with expert annotation than reproduced msiPL at the
        same budget?
        \emph{Justification:} the proposal's comparison (P3), made fair.
  \item \textbf{Ranking and representation.} Which component carries the
        result: the attribution method (IG versus direct correlation with the
        segmentation), or the nonlinear encoder (VAE versus PCA)?
        \emph{Justification:} without this ablation a gain cannot be credited
        to any component (\cref{sec:background}).
  \item \textbf{Spatial context.} Does neighbourhood information give a
        consistent, spatially specific gain when it enters the encoder, the
        segmentation or the ranking statistic?
        \emph{Justification:} P1 and P3, tested against shuffled-neighbour
        controls.
  \item \textbf{Reproducibility.} Do the msiPL and S3PL baselines reproduce
        their published scores, and how much do training seeds change the
        results?
        \emph{Justification:} baselines that do not reproduce cannot anchor a
        comparison.
\end{enumerate}

\subsection{Changes from the proposal}
\Cref{tab:proposal} lists each commitment in the proposal, what was
delivered, and why anything changed.

\begin{table*}[t]
\caption{Proposal commitments and outcomes.}
\label{tab:proposal}
\centering\small
\begin{tabularx}{\textwidth}{@{}p{3.6cm}X X@{}}
\toprule
Proposal commitment & Outcome & Reason for any change \\
\midrule
Neighbourhood input to a 1D VAE (P1) & Built: uniform $3\times3$ mean, learned attention, $5\times5$--$9\times9$ windows, centre masking & Delivered; extended with shuffled and zeroed controls to test spatial specificity \\
Spatial-coherence penalty on neighbouring latent means (P2) & Tested at $\lambda\in\{0,0.01,0.1,1\}$; $\lambda=0$ retained & Smoother latents but no gain in reconstruction, probe accuracy or Moran's~$I$ \\
Poisson noise augmentation, after DeepION (P2) & Tested; switched off & Did not pass its development gate \\
Bayesian optimisation of $\lambda$ in $[0.01, 1]$ & Omitted; a fixed grid was used & No penalty value passed the gate, so there was nothing to tune \\
Peak extraction from the L2 norm of first-layer weights & Kept as a comparator; replaced by GMM-targeted IG & First-layer weights ignore the nonlinear encoder; IG was later compared with direct correlation (RQ2) \\
Baselines: msiPL, S3PL, MALDIquant & msiPL and S3PL reproduced and audited; MALDIquant omitted; model-free rankings, random and supervised oracle added & Model-free rankings are a stronger and fairer reference than a spectrum-wise peak picker \\
Data: raw profile-mode imzML for GBM (PR001292) and CAC (MTBLS415) & Audited HDF5 conversions where needed; RCC added as a third, held-out dataset & Practical loading and a held-out confirmation set \\
mSCF1, PCC distributions, GPU memory and runtime & mSCF1 delivered with matched budgets, peak-level scoring and predeclared rules; memory and runtime measured; PCC distributions used in diagnostics & Bin-level scoring rewards duplicate bins of one peak \\
(Not in proposal) & Tests of the encoder (PCA), training seeds and S3PL neighbour topology & Added after reviews exposed missing controls \\
\bottomrule
\end{tabularx}
\end{table*}
